\section*{Hoofdstuk \ref{ch:g7:volume-mass-density} --- Volume, massa, dichtheid}

\begin{solution}{exo:g7:volume-mass-density:1}
$\rho = m/V$: $\rho$ de \omterm{def:g7:volume-mass-density:formula}{dichtheid}, $m$ de \omterm{def:g3:measuring-mass:mass}{massa} van het monster,
$V$ zijn \omterm{def:g6:mass-and-volume:volume}{volume}. Met \omterm{def:g3:measuring-mass:units}{gram} en \omterm{def:g6:mass-and-volume:volume}{kubieke centimeter} komt $\rho$
eruit in \unit{g/cm^3}.
\end{solution}

\begin{solution}{exo:g7:volume-mass-density:2}
$\rho = 270 \div 100 = \qty{2.7}{g/cm^3}$: aluminium.
\end{solution}

\begin{solution}{exo:g7:volume-mass-density:3}
$m = 9.0 \times 40 = \qty{360}{g}$ koper.
$V = 1930 \div 19.3 = \qty{100}{cm^3}$ goud --- het eigen getal van de
kroonzaak.
\end{solution}

\begin{solution}{exo:g7:volume-mass-density:4}
Tarreren haalt de eigen \omterm{def:g3:measuring-mass:mass}{massa} van de cilinder weg, zodat de balans de
\omterm{def:g3:solids-liquids-gases:liquid}{vloeistof} alleen meldt --- de eerlijke-nulregel. De formule krijgt dan
de \omterm{def:g3:measuring-mass:mass}{massa} $m$ van de balans en de meniscusaflezing $V$ te eten.
\end{solution}

\begin{solution}{exo:g7:volume-mass-density:5}
$\rho = 70 \div 50 = \qty{1.4}{g/cm^3}$ --- honingachtig. IJs met
$0.9$ is veel minder dicht: het \omterm{def:g1:floating-and-sinking:words}{drijft} er hoog in.
\end{solution}

\begin{solution}{exo:g7:volume-mass-density:6}
Beide eenheden schalen mee: een \omterm{def:g6:mass-and-volume:volume}{liter} is $1000$ \unit{cm^3} en een
\omterm{def:g3:measuring-mass:units}{kilogram} is $1000$ \omterm{def:g3:measuring-mass:units}{gram}, dus vallen de twee duizendtallen tegen
elkaar weg --- één getal dient allebei. \omterm{def:g3:measuring-mass:units}{Gram} met \omterm{def:g6:mass-and-volume:volume}{liter} mengen slaat
dat wegvallen over en levert een getal dat er duizendvoudig naast zit.
\end{solution}

\begin{solution}{exo:g7:volume-mass-density:7}
$\rho = 58 \div 5.0 = \qty{11.6}{g/cm^3}$ --- vlak bij de $11.3$ van
lood, nergens in de buurt van de $19.3$ van goud. Oordeel: geen zuiver
goud; waarschijnlijk een bedrieger met een loden hart in gouden kleren.
\end{solution}

\begin{solution}{exo:g7:volume-mass-density:8}
$100 \times 100 \times 100 = 1000000$ \unit{cm^3} --- een miljoen. Bij
\qty{1}{g} per stuk is dat \qty{1000000}{g} $=$ \qty{1000}{kg}: een
volle ton per kubieke \omterm{def:g2:measuring-length:units}{meter} water --- en daarom is een groot waterbed
meubilair voor de begane grond.
\end{solution}

\begin{solution}{exo:g7:volume-mass-density:9}
$V = 260 - 200 = \qty{60}{cm^3}$;
$\rho = 420 \div 60 = \qty{7.0}{g/cm^3}$ --- tussen aluminium ($2.7$)
en ijzer ($7.9$), het dichtst onder ijzer. Het tussenin liggen
fluistert: een mengsel of een verguld-verzilverde bedrieger, in elk
geval niet de $10.5$ van zilver.
\end{solution}

\begin{solution}{exo:g7:volume-mass-density:10}
\omterm{def:g3:measuring-mass:mass}{Massa} van de \omterm{def:g3:solids-liquids-gases:liquid}{vloeistof}: $1010 - 380 = \qty{630}{g}$;
\omterm{def:g6:mass-and-volume:volume}{volume}: $\qty{75}{cL} = \qty{750}{cm^3}$; dus
$\rho = 630 \div 750 = \qty{0.84}{g/cm^3}$ --- een olie.
\end{solution}

\begin{solution}{exo:g7:volume-mass-density:11}
\omterm{def:g7:volume-mass-density:formula}{Dichtheid} van het pakket: $1200 \div 4000 = \qty{0.3}{g/cm^3}$ ---
ver onder de $1.0$ van water: de boei \omterm{def:g1:floating-and-sinking:words}{drijft} hoog. De $2.7$ van
aluminium beschrijft alleen de metalen huid; het pakket bestaat vooral
uit ingesloten \omterm{def:g2:air-around-us:air}{lucht}, en het is het pakket --- totale \omterm{def:g3:measuring-mass:mass}{massa} over
totaal \omterm{def:g6:mass-and-volume:volume}{volume} --- dat tegenover de drijfregel staat.
\end{solution}

\begin{solution}{exo:g7:volume-mass-density:12}
Protocol: weeg de ketting ($m$); meet haar \omterm{def:g6:mass-and-volume:volume}{volume} met de
stijgmethode, helemaal ondergedompeld, zonder bellen ($V$); bereken
$\rho = m/V$; bij zuiverheid verwacht je ongeveer \qty{9.0}{g/cm^3}.
Eerlijke twijfels: holle schakels sluiten \omterm{def:g2:air-around-us:air}{lucht} op en blazen $V$ op,
wat een lage \omterm{def:g7:volume-mass-density:formula}{dichtheid} veinst; en een sluw mengsel (of een verguld
kernstuk) kan vlak bij $9.0$ landen zonder ook maar enig zuiver koper
te bevatten. Strakkere toets: herhaal de volumemeting nadat je de
schakels hebt volgezogen (schud de bellen eruit), en voeg een
onafhankelijke handtekening toe --- bijvoorbeeld de \omterm{def:g1:magnets:magnet}{magneet} (een
stalen kern verraadt zich onmiddellijk) of, in een werkplaats, een
gemeten smeltgedrag.
\end{solution}

\begin{solution}{pb:g7:volume-mass-density:1}
\textbf{1.} $5400 \div 2000 = \qty{2.7}{g/cm^3}$: aluminium.
\textbf{2.} $1800 \div 200 = \qty{9.0}{g/cm^3}$: koper.
\textbf{3.} $4520 \div 400 = \qty{11.3}{g/cm^3}$: lood --- dicht,
zacht en giftig bij achteloos hanteren: handschoenen.
\textbf{4.} $8100 \div 1500 = \qty{5.4}{g/cm^3}$ --- tussen aluminium
($2.7$) en ijzer ($7.9$): een gemengde zak van die twee, qua
\omterm{def:g6:mass-and-volume:volume}{volume} ruwweg half om half.
\textbf{5.} $V = 540 \div 2.7 = \qty{200}{cm^3}$.
\textbf{6.} Eén kubus: $m = 9.0 \times 25 = \qty{225}{g}$. De
\qty{1800}{g} van de rol leveren $1800 \div 225 = 8$ volle kubussen.
\textbf{7.} $\qty{20}{kg} = \qty{20000}{g}$;
$20000 \div 225 = 88.9$: $88$ kubussen.
\textbf{8.} $m = 7.9 \times 3000 = \qty{23700}{g} \approx
\qty{23.7}{kg}$.
\textbf{9.} $3860 \div 200 = \qty{19.3}{g/cm^3}$ --- een perfecte
match met goud.
\textbf{10.} Wolfraam --- \omterm{def:g7:volume-mass-density:formula}{dichtheid} $19.3$, dezelfde handtekening tot
op de decimaal. De dichtheidstoets wijst kandidaten aan; ze kan
stoffen die toevallig een handtekening delen niet uit elkaar houden.
\textbf{11.} Geen van beide instrumenten heeft gefaald: \omterm{def:g3:measuring-mass:mass}{massa} en
\omterm{def:g6:mass-and-volume:volume}{volume} kloppen, en de deling ook. De toets heeft eerlijk vastgesteld
dat de staaf \emph{ofwel} goud is \emph{ofwel} iets met precies de
\omterm{def:g7:volume-mass-density:formula}{dichtheid} van goud --- ze heeft de verdachten tot twee teruggebracht.
\textbf{12.} Verdedigbare keuzes: de klank en het gevoel van een
aangeslagen staaf (wolfraam is veel harder --- een vijl- of
hardheidstoets op een verborgen hoekje scheidt ze snel), of een erkende
deskundige meting van hoe de staaf \omterm{def:g5:heat-and-insulation:heat}{warmte} of stroom geleidt (goud
behoort tot de beste \omterm{def:g4:conductors-insulators:conductor}{geleiders}, wolfraam blijft ver achter) ---
zachter dan welke oven ook, en samen met \omterm{def:g7:volume-mass-density:formula}{dichtheid} beslissend.
\textbf{13.} Bijvoorbeeld: ``\omterm{def:g7:volume-mass-density:formula}{Dichtheid} kan bewijzen dat een stof
\emph{niet} is wat ze beweert. Ze kan alleen suggereren wát ze is ---
handtekeningen beperken de verdachten. En in haar eentje kan ze nooit
veroordelen: identiteit wil twee onafhankelijke getuigen.''
\end{solution}
